\section{CONCLUSION}

As spatial computing and generative artificial intelligence converge, the necessity for robust, secure, and highly optimized integration architectures becomes paramount. In this research, the Vectra Spatial Computing Protocol was designed, implemented, and rigorously evaluated to bridge the gap between high-fidelity digital twins and localized generative AI pipelines. 

It was demonstrated that computationally expensive deep learning models—specifically U\textsuperscript{2}-Net, SDXL-Lightning, and TripoSR—can be successfully orchestrated on constrained, consumer-grade edge hardware (strictly within an 8GB VRAM threshold). By enforcing a decoupled, asynchronous client-server architecture, the visual presentation layer was entirely isolated from the heavy tensor operations of the inference backend, ensuring that real-time spatial navigability was never compromised by algorithmic latency. 

Furthermore, the introduction of the Deep Splat Excavation (DBSE) algorithm resolved a critical spatial occlusion problem inherent to dense Gaussian Splatting environments. Through non-destructive, shader-level volumetric masking, it was proven that synthetic 3D assets can be surgically injected into digitized physical spaces without destructive clipping, Z-fighting, or permanent alteration of the foundational point cloud data. 

Ultimately, this research establishes a vital baseline for the future of localized spatial computing. As autonomous agents and generative models are increasingly granted spatial agency within digital twins, relying solely on heuristic software boundaries will prove insufficient. The architectural solutions presented herein lay the groundwork for embedding definitive mathematical safeguards, such as Control Barrier Functions, directly into the spatial rendering pipeline. It is concluded that securing the intersection of generative AI and spatial computing at the edge hardware level is not merely an optimization challenge, but a fundamental prerequisite for the safe and reliable deployment of future digital infrastructures.